\documentclass[10pt,a4paper]{amsart}
\usepackage[margin=19mm]{geometry}
\usepackage{amsmath,amssymb,mathtools}
\usepackage[scr=rsfs,cal=euler]{mathalfa}
\usepackage{xfrac}
\usepackage[unicode,hidelinks,bookmarks=false]{hyperref}
\newcommand{\ie}{i.e.\ }
\newcommand{\cf}{cf.\ }
\newcommand{\resp}{respectively\ }
\newcommand{\Subsection}{\subsection}
\providecommand{\nobreakdash}{\nobreak\hbox{-}}
\setlength{\emergencystretch}{2em}
\multlinegap0pt
\usepackage{bm}
\usepackage{algorithm}\usepackage{algpseudocode}
\usepackage[shortlabels]{enumitem}
\usepackage{multirow}
\usepackage{booktabs}
\usepackage{mathtools}
\usepackage{dsfont}
\usepackage{tikz}
\newtheorem{lemma}{Lemma}
\newcommand{\bb}{{\bf b}}
\newcommand{\bbN}{{\mathbb N}}
\newcommand{\bbR}{{\mathbb R}}
\newcommand{\bd}{{\bf d}}
\newcommand{\be}{\begin{equation}}
\newcommand{\bean}{\setlength\arraycolsep{1pt}\begin{eqnarray*}}
\newcommand{\bx}{{\bf x}}
\def\cF{\mathcal F}
\newcommand{\ee}{\end{equation}}
\newcommand{\eean}{\end{eqnarray*}}
\title{Nonparametric Estimation of a Factorizable Density using Diffusion Models}
\author{Hyeok Kyu Kwon, Dongha Kim, Ilsang Ohn, Minwoo Chae}
\date{Source: arXiv:2501.01783v3 (CC BY 4.0)}
\begin{document}
\maketitle

\noindent\emph{Source and licence.} Nonparametric Estimation of a Factorizable Density using Diffusion Models. Version arXiv:2501.01783v3 (CC BY 4.0), distributed by arXiv under the Creative Commons Attribution 4.0 International licence, http://creativecommons.org/licenses/by/4.0/, as recorded in the arXiv licence metadata for this exact version and shown on the abstract page https://arxiv.org/abs/2501.01783v3. Full licence notice: \texttt{licenses/arxiv-2501.01783-CC-BY-4.0.txt}.\par\smallskip

\noindent\textit{Editorial scope.} Counted unit: the article's lemma secnn:exp (source0.tex lines 1866--1879). The article's complete original proof of it is delivered with it (source0.tex lines 1881--1998, one \texttt{proof} environment of 118 lines). No mathematics is written, repaired or re-derived: the statement and the proof are the article's own lines, in the article's own notation, and no retained line is substituted or re-flowed.  Retained uncounted as the source-local context the proof needs: lines 1484--1491; lines 1497--1518; lines 1524--1531; lines 1687--1700; lines 1705--1716.  Nothing was omitted from any retained span, so the package carries no omission record.  No retained line contains a citation, so the wrapper carries no bibliography and no bibliography entry.  No retained line contains a figure environment, an image-inclusion command or drawing code, so no image is delivered and the package declares no asset dependency.  Editorial formatting only, in addition: an AMS \texttt{amsart} preamble that restates the article's own declarations and macros used by the retained lines, namely the environments \texttt{lemma} on separate counters, together with the standard packages those lines need.  The retained environments print in this document as Lemma 1 in order of appearance, whereas the article prints the same items with the numbering of its own full document.  The complete native source of the article as distributed in the arXiv e-print for this exact version is bundled unchanged as \texttt{sources/171117/source0.tex}.
\begin{lemma}[Concatenation]
\label{secnn:comp}
Let $K \in \bbN, \{d_1,\ldots,d_{K+1}\} \subseteq \bbN$ be given. Consider $L^{(k)} \in \bbN_{\geq 2}$, $ s^{(k)}, M^{(k)} > 0$ and $\bd^{(k)} = (d_1^{(k)},\ldots,d_{L^{(k)}}^{(k)})^{\top} \in \bbN^{L^{(k)}}$ with $d_1^{(k)} = d_{k}$ and $ d_{L^{(k)}}^{(k)} = d_{k+1}$ for $k \in [K]$. For any neural networks $f_1,\ldots,f_{K}$ with $f_k \in \cF_{\rm NN}(L^{(k)},\bd^{(k)},s^{(k)}, M^{(k)}), k \in [K]$, there exists a neural network $f \in \cF_{\rm NN}(L,\bd,s,M)$ with 
\bean
L = \sum_{k=1}^{K} L^{(k)}, \quad \| \bd \|_{\infty} \leq 2 \max_{k \in [K]} \|\bd^{(k)}\|_{\infty}, \quad s \leq 2 \sum_{k=1}^{K} s^{(k)}, \quad M = \max_{k \in [K]} M^{(k)}
\eean
such that $f(\bx) = (f_K \circ \cdots \circ f_{1})(\bx) $ for any $\bx \in \bbR^{d_1}$.
\end{lemma}

\begin{lemma}[Parallelization]
\label{secnn:par}
Let $K \in \bbN$ be given. Consider $L^{(k)} \in \bbN_{\geq 2}$, $ s^{(k)}, M^{(k)} > 0$, $\bd^{(k)} = (d_1^{(k)},\ldots,d_{L^{(k)}}^{(k)})^{\top} \in \bbN^{L^{(k)}}$ for $k \in [K]$.
For any neural networks $f_1,\ldots,f_{K}$ with 
\bean
f_k \in \cF_{\rm NN} (L^{(k)}, \bd^{(k)}, s^{(k)}, M^{(k)} ), \quad k \in [K],
\eean
there exists a neural network $f \in \cF_{\rm NN}(L,\bd,s,M)$ with 
\bean
&& L = \max_{k \in [K]} L^{(k)}, \quad \|\bd\|_{\infty} \leq 2 \sum_{k=1}^{K} \|\bd^{(k)}\|_{\infty},
\\
&& s \leq 2\sum_{k=1}^{K} \left(s^{(k)} + L d_{L^{(k)}}^{(k)} \right), \quad M \leq \left(\max_{k \in [K]} M^{(k)} \right) \vee 1
\eean
such that
\bean
f(\bx) = \left( f_1(\bx^{(1)}),\ldots, f_K(\bx^{(K)}) \right) \in \bbR^{d_{L^{(1)}}^{(1)} + \cdots + d_{L^{(K)}}^{(K)}}
\eean
for $\bx = (\bx^{(1)},\ldots,\bx^{(K)}) \in \bbR^{d_1^{(1)}+\cdots+d_{1}^{(K)} }$. If $L^{(1)} = \cdots = L^{(K)} = \widetilde L$ with $\widetilde L \in \bbN_{\geq 2}$, $(L,\bd,s,M)$ also satisfies
\bean
L = \widetilde L, \quad \|\bd\|_{\infty} \leq \sum_{k=1}^{K} \|\bd^{(k)}\|_{\infty}, \quad s \leq \sum_{k=1}^{(K)} s^{(k)}, \quad M \leq  \max_{k \in [K]} M^{(k)}.
\eean
\end{lemma}

\begin{lemma}[Linear function]
\label{secnn:lin}
Let $W \in \bbR^{d_2 \times d_1}, \bb \in \bbR^{d_2}$ be given with $d_1,d_2 \in \bbN$. There exists a neural network $f_{\rm lin} \in \cF_{\rm NN}(L,\bd,s,M)$ with
\bean
L = 2, \quad \bd = (d_1,2d_2,d_2)^{\top}, \quad s = 2\|W\|_{0} + 2\|\bb\|_{0} + 2d_2, \quad M = \max\{ \|W\|_{\infty}, \|\bb\|_{\infty}, 1\}
\eean
such that $f_{\rm lin}(\bx) = W\bx + \bb$ for any $\bx \in \bbR^{d_1}$.
\end{lemma}

\begin{lemma}[Multiplication]
\label{secnn:mult}
Let $m \geq 2, C \geq 1, 0 < \widetilde \epsilon \leq 1$ be given. For any $\epsilon > 0$, there exists a positive constant $C_{N,1}$ and a neural network $f_{\rm mult} \in \cF_{\rm NN}(L,\bd,s,M)$ with 
\bean
&& L \leq C_{N,1} \log m \{ \log (1/\epsilon) + m \log C \}, \quad \bd = (m,48m,\ldots,48m,1)^{\top},
\\
&& s \leq C_{N,1} m \{ \log(1/\epsilon) +  \log C \}, \quad M = C^{m}
\eean
such that
\bean
\left\vert f_{\rm mult}(\widetilde \bx) - \prod_{i=1}^{m} x_{i} \right\vert \leq \epsilon + mC^{m-1}\widetilde \epsilon, \quad \forall \bx \in [-C,C]^{m}, \widetilde \bx \in \bbR^{m} \ {\rm with} \ \|\bx-\widetilde \bx\|_{\infty} \leq \widetilde \epsilon,
\eean
$ \| f_{\rm mult} \|_{\infty} \leq C^{m}$ and $f_{\rm mult}(\widetilde \bx) = 0$ if $0 \in \{\widetilde x_1,\ldots,\widetilde x_{m}\}$.
\end{lemma}

\begin{lemma}[Clipping function]
\label{secnn:clip}
Let $\underline{\bb} = (\underline{b}_1,\ldots,\underline{b}_m), \overline{\bb} = (\overline{b}_1,\ldots,\overline{b}_m) \in \bbR^{m}$ be given with $m \in \bbN$ and $\underline{b}_i \leq \overline{b}_i$ for all $i \in [m]$. Then, there exists a neural network 
\bean
f_{\rm clip}^{(\underline{\bb},\overline{\bb})} \in \cF_{\rm NN}(2,(m,2m,m)^{\top}, 7m, \| \underline{\bb} \|_{\infty} \vee \| \overline{\bb} \|_{\infty} )
\eean
such that 
\bean
f_{\rm clip}^{(\underline{\bb},\overline{\bb})} (\bx) = \left( \overline{b}_1 \wedge \{x_1 \vee \underline{b}_1 \}, \ldots, \overline{b}_m \wedge \{ x_m \vee \underline{b}_m \}  \right) \in \bbR^{m}
\eean
for $\bx = (x_1,\ldots,x_m) \in \bbR^{m}$.
\end{lemma}

\begin{lemma}[Negative exponential function]
\label{secnn:exp}
For any $0 < \epsilon < 2^{-4e+2}$, there exists a positive constant $C_{N,3}$ and a neural network $f_{\exp} \in \cF_{\rm NN}(L,\bd,s,M)$ with
\bean
&& L \leq C_{N,3} \log(1/\epsilon) \log \log (1/\epsilon), \quad \|\bd\|_{\infty} \leq C_{N,3}  \{ \log(1/\epsilon)\}^{3}
\\
&& s \leq C_{N,3}  \{ \log(1/\epsilon)\}^{4}, \quad M \leq C_{N,3} \epsilon^{-1}
\eean
such that
\bean
\left\vert e^{-x} - f_{\exp}(\widetilde x) \right\vert \leq \epsilon + \vert x-\widetilde x\vert
\eean
for any $x \geq 0$ and $\widetilde x \in \bbR$.
\end{lemma}

\begin{proof}
Let $0 < \epsilon < 2^{-4e+2}, D_1 = \lfloor \log (4/\epsilon) \rfloor +1, D_2 = \lfloor \log (4/\epsilon) / \log 2  \rfloor +1$ and $\underline{T}_i = i-1, \overline{T}_i = i+1$ for $i \in [D_1]$. Then, Taylor's theorem yields that for any $i \in [D_1]$ and $x \in [\underline{T}_i, \overline{T}_i]$,
\bean
e^{-x} = e^{-\underline{T}_i} e^{-(x-\underline{T}_i)} = e^{-\underline{T}_i} \left\{ P_i(x)
 + \frac{(-1)^{D_2} e^{-\xi(x-\underline{T}_i)} (x-\underline{T}_i)^{D_2} }{D_2 !}
\right\}
\eean
for a suitable $\xi \in [0,1]$, where
\bean
P_i(x) = 1 - (x-\underline{T}_i) + \sum_{k=2}^{D_2-1} \frac{(-1)^{k} (x-\underline{T}_i)^{k}} {k!}.
\eean
Since $0 \leq x-\underline{T}_i \leq 2$ and $\underline{T}_i \geq 0$, it follows that
\be
\left\vert e^{-x} - e^{-\underline{T}_i} P_i(x) \right\vert  \leq \frac{2^{D_2}}{D_2 !} \leq \left(\frac{2e}{D_2} \right)^{D_2} \leq \left(\frac{1}{2} \right)^{D_2} \leq \frac{\epsilon}{4}, \label{eq:expremainder}
\ee
where the second inequality holds because $k! \geq k^{k} e^{-k}$ for any $k \in \bbN$.
Let $N_1$ be a constant in Lemma~\ref{secnn:mult}. For $k \geq 2$, there exists a neural network $f_{\rm mult}^{(k)} \in \cF_{\rm NN}(L_{\rm mult}^{(k)}, \bd_{\rm mult}^{(k)}, s_{\rm mult}^{(k)}, M_{\rm mult}^{(k)})$ with 
\be \begin{split}
& L_{\rm mult}^{(k)} \leq N_1  \log k \{ \log (4D_2/\epsilon^2) +  k \log 2 \}, \quad \bd_{\rm mult}^{(k)} = (k,48k,\ldots,48k,1)^{\top}, \quad
\\
& s_{\rm mult}^{(k)} \leq N_1 k \{ \log(4D_2/\epsilon^2) + \log 2 \}, \quad M_{\rm mult}^{(k)} = 2^{k} \label{eq:nnmultexpcomp}
\end{split} \ee
such that $\vert f_{\rm mult}^{(k)}(x_1,\ldots,x_k) - \prod_{i=1}^{k} x_i \vert \leq \epsilon^2/(4D_2)$ for any $x_1,\ldots,x_k \in [-2,2]$.
For any $k \geq 1$ and $i \in [D_1]$, Lemma~\ref{secnn:lin} implies that there exists a neural network 
\bean
f_{\rm lin}^{(i,k)} \in \cF_{\rm NN}(2,(1,2k,k)^{\top}, 6k, \underline{T}_i)
\eean
such that $f_{\rm lin}^{(i,k)}(x) = (x-\underline{T}_i,\ldots,x-\underline{T}_i)^{\top} \in \bbR^{k}$ for any $x \in \bbR$. 
Combining Lemma~\ref{secnn:comp} with the last display, it follows that $f_{\rm pow}^{(i,k)} = f_{\rm mult}^{(k)} \circ f_{\rm lin}^{(i,k)} \in \cF_{\rm NN}(L_{\rm pow}^{(i,k)},\bd_{\rm pow}^{(i,k)},s_{\rm pow}^{(i,k)}, M_{\rm pow}^{(i,k)})$ for $i \in [D_{1}], k \geq 2$ with
\bean
L_{\rm pow}^{(i,k)} =  L_{\rm mult}^{(k)} +2, \quad \| \bd_{\rm pow}^{(i,k)}\|_{\infty} \leq 96k, \quad s_{\rm pow}^{(i,k)} = 2s_{\rm mult}^{(k)} + 12k, \quad M_{\rm pow}^{(i,k)} = \underline{T}_i \vee M_{\rm mult}^{(k)}
\eean
and
\bean
\left\vert f_{\rm pow}^{(i,k)}(x) - (x-\underline{T}_i)^{k} \right\vert \leq \frac{ \epsilon^2}{4D_2}
\eean
for $x \in [\underline{T}_i,\overline{T}_i]$.
Consider functions $f_1,\ldots,f_{D_1} : \bbR \to \bbR$ such that
\bean
f_i(\cdot) = 1 - f_{\rm lin}^{(i,1)}(\cdot) +\sum_{k=2}^{D_2-1} \frac{(-1)^{k}   f_{\rm pow}^{(i,k)}(\cdot)  }{k!}, \quad i \in [D_1].
\eean
Since $f_i - 1$ is a linear combination of $ f_{\rm lin}^{(i,1)}, f_{\rm pow}^{(i,2)}, \ldots,f_{\rm pow}^{(i,D_2-1)}$, Lemma~\ref{secnn:comp}, Lemma~\ref{secnn:par} and Lemma~\ref{secnn:lin} implies that $f_i \in \cF_{\rm NN}(L^{(i)},\bd^{(i)},s^{(i)},M^{(i)})$ with
\be \begin{split}
& L^{(i)} \leq L_{\rm pow}^{(i,D_2-1)} + 2 \leq D_3 \log D_2  \{ \log (1/\epsilon) + D_2 \} 
\\
& \| \bd^{(i)} \|_{\infty} \leq 2\max \left(2\sum_{k=2}^{D_2-1} \| \bd_{\rm pow}^{(i,k)}\|_{\infty} + 4 , D_2-1 \right) \leq D_3 D_2^2
\\
& s^{(i)} \leq 2\left\{ \sum_{k=2}^{D_2-1} \left( s_{\rm pow}^{(i,k)} +L_{\rm pow}^{(i,D_2-1)} + 2 \right) + L_{\rm pow}^{(i,D_2-1)} + 8 \right\} + 2D_2 + 2
\\
& \leq D_3 D_2^{2} \{ \log (1/\epsilon) + D_2 \} 
\\
& M^{(i)} \leq \left( \max_{k \in [D_2-1]}  M_{\rm pow}^{(i,k)} \right) \vee 1 \leq (D_1+1) \vee 2^{D_2-1} \label{eq:nnexpcomp}
\end{split} \ee
for a large enough constant $D_3 = D_3(N_1)$.
Then, 
\bean
&& \left\vert P_i(x) - f_{i}(x) \right\vert \leq \sum_{k=2}^{D_2-1} \frac{\vert f_{\rm pow}^{(i,k)}(x) - (x-\underline{T}_i)^{k}  \vert}{k!}
\\
&& \leq D_2 \max_{2 \leq k \leq D_2-1}\left\vert f_{\rm pow}^{(i,k)}(x) - (x-\underline{T}_i)^{k} \right\vert \leq \frac{\epsilon^2}{4}, \quad i \in [D_{1}]
\eean
for $x \in [\underline{T}_i,\overline{T}_i]$.
Combining (\ref{eq:expremainder}) with the last display, we have
\be
\left\vert e^{-x} - e^{-\underline{T}_i} f_i(x) \right\vert \leq \left\vert e^{-x} - e^{-\underline{T}_i} P_i(x) \right\vert + \left\vert P_i(x) - f_{i}(x) \right\vert \leq \frac{\epsilon}{4}+ \frac{\epsilon^2}{4} \leq \frac{\epsilon}{2}, \quad i \in [D_1] \label{eq:nnexp2}
\ee
for $x \in [\underline{T}_i,\overline{T}_i]$, where the first inequality holds because $e^{-\underline{T}_i} \leq 1$.
Consider functions $f_{\rm swit}^{(1)},\ldots,f_{\rm swit}^{(D_1+1)} : \bbR \to [0,1]$ such that
\bean
&& f_{\rm swit}^{(1)}(\cdot) = \frac{1}{\overline{T}_{1}-\underline{T}_{2}} \rho \left(-f_{\rm clip}^{(\underline{T}_{2},\overline{T}_{1})} (\cdot) + \overline{T}_{1} \right),
\\
&& f_{\rm swit}^{(i)}(\cdot) = \frac{1}{\overline{T}_{i-1}-\underline{T}_{i}} \rho \left(f_{\rm clip}^{(\underline{T}_
{i},\overline{T}_{i-1})} (\cdot) - \underline{T}_{i} \right) - \frac{1}{\overline{T}_{i}-\underline{T}_{i+1}} \rho \left(f_{\rm clip}^{(\underline{T}_{i+1},\overline{T}_{i})} (\cdot) - \underline{T}_{i+1} \right),
\\
&& 2 \leq i \leq D_1,
\\
&& f_{\rm swit}^{(D_1+1)}(\cdot) = \frac{1}{\overline{T}_{D_1}-\underline{T}_{D_1+1}} \rho \left(f_{\rm clip}^{(\underline{T}_
{D_1+1},\overline{T}_{D_1})} (\cdot) - \underline{T}_{D_1+1} \right),
\eean
where $\underline{T}_{D_{1}+1} = D_1$ and $f_{\rm clip}^{(\underline{T}_i,\overline{T}_{i-1})} \in \cF_{\rm NN}(2, (1,2,1)^{\top}, 7, D_1)$ denotes the neural network in Lemma~\ref{secnn:clip}.
Note that $\sum_{i=1}^{D_1+1} f_{\rm swit}^{(i)}(x) = 1$ for $x \in \bbR$, and $f_{\rm swit}^{(i)}(x) = 0$ for $x \in [0,\underline{T}_i] \cup [\overline{T}_i, \infty), i \in [D_1]$ and $f_{\rm swit}^{(D_1+1)}(x) = 0$ for $x \leq \underline{T}_{D_1+1}$.
Consider a function $f: \bbR \to \bbR$ such that $f(\cdot) = \sum_{i=1}^{D_1}  e^{-\underline{T}_i}  f_{\rm mult}^{(2)}(f_{\rm swit}^{(i)}(\cdot) , f_{i}(\cdot))$.
Since $\vert f_{\rm mult}^{(2)} (x_1,x_2) - x_1x_2\vert \leq \epsilon^2/(4D_2)$ for any $x_1,x_2 \in [-2,2]$, we have
\bean
&& \left\vert e^{-x} - f(x) \right\vert \leq \left\vert e^{-x} - \sum_{i=1}^{D_1} e^{-\underline{T}_i} f_{\rm swit}^{(i)}(x) f_i(x) \right\vert + \frac{D_1\epsilon^{2}}{4D_2}  
\\
&& = \left\vert  \sum_{i=1}^{D_1} f_{\rm swit}^{(i)}(x) \left\{ e^{-x} -  e^{-\underline{T}_i}f_{i}(x) \right\} + f_{\rm swit}^{(D_1+1)}(x)e^{-x} \right\vert + \frac{D_1\epsilon^{2}}{4D_2},
\\
&& \leq \sum_{i=1}^{D_1} f_{\rm swit}^{(i)}(x) \left\vert e^{-x} - e^{-\underline{T}_i}f_{i}(x) \right\vert + \frac{\epsilon}{4} + \frac{D_1\epsilon^{2}}{4D_2} 
\\
&& \leq \frac{\epsilon}{2} \sum_{i=1}^{D_1} f_{\rm swit}^{(i)}(x) + \frac{\epsilon}{4} + \frac{D_1\epsilon^{2}}{4D_2}
\leq \left(\frac{3}{4} + \frac{(D_1+1)\epsilon}{4D_2} \right) \epsilon  \leq \epsilon
\eean
for $x \in [0,\infty)$, where the second inequality holds because  $\vert f_{\rm swit}^{(D_1+1)}(x)e^{-x} \vert \leq e^{-x} \leq \epsilon/4$ for $x \geq \underline{T}_{D_1+1}$ and the third inequality holds by  (\ref{eq:nnexp2}).
Combining (\ref{eq:nnmultexpcomp}) and (\ref{eq:nnexpcomp}) with Lemma~\ref{secnn:par} and Lemma~\ref{secnn:comp}, we have $  f_{\rm mult}^{(2)}(f_{\rm swit}^{(i)}(\cdot) , f_{i}(\cdot)) \in \cF_{\rm NN}(\widetilde L^{(i)},\widetilde \bd^{(i)},\widetilde s^{(i)}, \widetilde M^{(i)})$ for $i \in [D_{1}]$ with
\bean
&& \widetilde L^{(i)} \leq (L^{(i)} \vee 2) + L_{\rm mult}^{(2)} \leq D_4 \log D_2  \{ \log (1/\epsilon) + \log D_2 \} 
\\
&& \| \widetilde \bd^{(i)} \|_{\infty} \leq 2\max \left( 2\|\bd^{(i)}\|_{\infty} + 4 ,  \|\bd_{\rm mult}^{(2)}\|_{\infty}   \right) \leq D_4 D_2^2
\\
&& \widetilde s^{(i)} \leq 4 s^{(i)} + 4 (L^{(i)} \vee 2) + 2 s_{\rm mult}^{(2)} + 28  \leq D_4 D_2^{2}  \{ \log (1/\epsilon) +  D_2 \} 
\\
&& \widetilde M^{(i)} \leq \max\left(D_1, M^{(i)}, M_{\rm mult}^{(2)}, 1 \right),
\eean
where $D_4 = D_4(D_3)$ is a large enough constant.
Since $f$ is a linear combination of $f_{\rm mult}^{(2)}(f_{\rm swit}^{(i)}(\cdot) , f_{i}(\cdot))$ for each $i \in [D_1]$, Lemma~\ref{secnn:comp}, Lemma~\ref{secnn:par} and Lemma~\ref{secnn:lin} implies that $f \circ \rho_{1} \in \cF_{\rm NN}(L,\bd,s,M)$ with
\bean
&& L \leq \max_{i \in [D_1]} \widetilde L^{(i)} + 3 \leq D_5 \log(1/\epsilon) \log \log (1/\epsilon)
\\
&& \|\bd\|_{\infty} \leq 2\max\left( 2\sum_{i=1}^{D_1} \| \widetilde \bd^{(i)}\|_{\infty}, 4 D_1 \right) \leq D_5  \{ \log(1/\epsilon)\}^{3}
\\
&& s \leq 2\sum_{i=1}^{D_1}\left(\widetilde s^{(i)} +  \max_{j \in [D_1]} \widetilde L^{(j)} \right) + 4D_1 + 4 \leq D_5  \{ \log(1/\epsilon)\}^{4}
\\
&& M \leq \max_{i \in [D_1]} \left(\widetilde M^{(i)} \vee e^{-\underline{T}_i} \right) \leq D_5 \epsilon^{-1}
\eean
for large enough constant $D_5 = D_5(D_4)$.
Note that $\vert(f \circ \rho_{1})(\widetilde x) - e^{-x} \vert \leq \vert(f \circ \rho_{1})(\widetilde x) - e^{-(\widetilde x \vee 0)} \vert \ + \vert e^{-(\widetilde x \vee 0)} - e^{-x} \vert \leq \epsilon + \vert x-\widetilde x\vert$ for any $x \geq 0$ and $\widetilde x \in \bbR$.
Then, the assertion follows by re-defining the constant.
\end{proof}

\end{document}
